\newcommand{\podstawy}{%
\par\smallskip
Na jednorodnym papierze przewodzącym stacjonarny przepływ prądu
modeluje pole elektrostatyczne w obszarze bez ładunków, jeśli odpowiadają
sobie warunki brzegowe. Potencjał spełnia równanie Laplace'a.
\[
\boxed{\nabla^2V=0,\qquad \mathbf E=-\nabla V,\qquad
|\mathbf E|=\sqrt{E_x^2+E_y^2}.}
\]
Linie pola są prostopadłe do ekwipotencjalnych i skierowane od wyższego
potencjału do niższego. Wartość natężenia i składowa ze znakiem to różne wielkości.
}
\newcommand{\teoriaA}{%
\par\smallskip
Dla idealnych równoległych elektrod przy $V(0)=0$ i $V(d)=U$:
\[
\boxed{V(x)=\frac Ud x,\qquad E_x=-\frac Ud,\qquad |\mathbf E|=\frac Ud.}
\]
Zależność opisuje wnętrze z pominięciem efektów brzegowych;
poza elektrodami pole jest niejednorodne.
}
\newcommand{\teoriaB}{%
\par\smallskip
Dla elektrody wewnętrznej o promieniu $r_w$ i potencjale $U$
oraz zewnętrznej $r_z$ o potencjale 0:
\[
\boxed{V(r)=U\frac{\ln(r_z/r)}{\ln(r_z/r_w)},\qquad
E_r(r)=\frac{U}{r\ln(r_z/r_w)}.}
\]
Wtedy $V(r_w)=U$, $V(r_z)=0$, a pole jest skierowane na zewnątrz.
Odwrócenie biegunowości odwraca wektor. Zmiana zera potencjału
o stałą nie zmienia pola; nie jest tym samym co zmiana znaku całego $V$.
}
\newcommand{\teoriaC}{%
\par\smallskip
Dla dowolnego układu elektrod nie zakładamy postaci liniowej ani
radialnej. Linie wyznaczamy z pomiarów potencjału i warunków na elektrodach.
Obszary bez danych oznaczamy, nie wypełniając ich fikcyjnymi pomiarami.
}
\newcommand{\roznice}{%
\par\smallskip
Dla trzech kierunków w tej samej odległości:
\[
\boxed{\bar V_i=\frac{V_{a,i}+V_{b,i}+V_{c,i}}3,\quad
s_i^*=\frac{s_i+s_{i+1}}2,\quad
E_{s,i}=-\frac{\bar V_{i+1}-\bar V_i}{s_{i+1}-s_i}.}
\]
Współrzędna $s$ oznacza $x$ lub $r$. W mianowniku używaj metrów:
$1\ \mathrm{mm}=10^{-3}\ \mathrm m$. Przypisz pochodną do środka przedziału.
Trzy kierunki nie muszą być trzema niezależnymi powtórzeniami tego samego pomiaru.
}
\newcommand{\mapowanie}{%
Linie ekwipotencjalne rysuj co 1--2 V. Szukane położenie między
zmierzonymi punktami wyznacz przez interpolację, nie przez zastąpienie
pomiarów idealnym polem. Dla sąsiadów odległych o $h,k$:
\[
\boxed{E_x\approx-\frac{V(x+h,y)-V(x,y)}h,\qquad
E_y\approx-\frac{V(x,y+k)-V(x,y)}k.}
\]
Powyższe składowe są różnicami w przód przy punkcie bazowym;
podaj użyte punkty i kroki. Dla wektora w środku komórki użyj
uśrednionych różnic na przeciwległych krawędziach. Nie łącz składowych
liczonych w różnych miejscach bez wyjaśnienia.
}
\newcommand{\dokladnosc}{%
Zapisz specyfikację i zakres woltomierza, krok siatki oraz stabilność
zasilania i kontaktu sondy. Dla niezależnych odczytów i odstępu $\Delta s$:
\[
\boxed{u(E_s)\approx\sqrt{
\frac{u^2(V_i)+u^2(V_{i+1})}{(\Delta s)^2}
+\frac{(V_{i+1}-V_i)^2}{(\Delta s)^4}u^2(\Delta s)}.}
\]
Uwzględnij korelacje, jeśli odczyty mają wspólne odniesienie lub błąd skali.
Bez danych przyrządów nie wpisuj wymyślonych słupków niepewności.
Różniczkowanie nasila wpływ szumu; interpolacja nie zwiększa liczby pomiarów.
}
\newcommand{\parametry}[1]{%
\begin{tabularx}{\linewidth}{|>{\raggedright\arraybackslash}p{55mm}|Y|}\hline
\naglowek\textbf{Parametr wariantu #1} & \textbf{Wpis ze stanowiska}\\\hline
Wariant i oznaczenie płyty &\rule{0pt}{9mm}\\\hline
Napięcia obu elektrod / odniesienie &\rule{0pt}{9mm}\\\hline
Geometria: $d$ albo $r_w,r_z$ &\rule{0pt}{9mm}\\\hline
Krok otworów $h,k$, początek i zwrot osi &\rule{0pt}{9mm}\\\hline
Kierunki a,b,c / kąty promieni &\rule{0pt}{9mm}\\\hline
Woltomierz: model, zakres, dokładność &\rule{0pt}{9mm}\\\hline
\end{tabularx}\par\medskip
Napięcie zalecane w instrukcji: 10 V. Zapisz rzeczywistą wartość
i biegunowość; nie przenoś wymiarów z cudzego przykładu.
}
\newcommand{\aparatura}{%
\par\medskip\noindent
\begin{minipage}[c]{.49\linewidth}\centering
\includegraphics[width=\linewidth]{../../Wspolne/Rysunki/uklad.pdf}
\end{minipage}\hfill
\begin{minipage}[c]{.47\linewidth}
\raggedright
\captionof{figure}{Układ pomiarowy z instrukcji 31\_wykon.pdf, s.~2.}
\smallskip Zasilacz wymusza różnicę potencjałów elektrod.
Woltomierz z sondą mierzy napięcie wybranego punktu względem
ustalonej elektrody odniesienia. Papier przewodzący stanowi model pola.
Wymiana płyty zmienia geometrię elektrod, nie zasadę pomiaru.
\end{minipage}\par\medskip}
